\chapter{Of Citizenship and the Census}
\label{art:citizenship}

\articlenote{In which is set down how a citizen is named, how a citizen comes into full standing, and how the silo keeps account of its own numbers, that none may be forgotten and none counted twice.}

\section{The Standing of a Citizen}

Citizenship, as Article 1 sets down, belongs to every person born within the silo and to every person who has taken the citizen's oath before Judicial upon reaching majority. This section and those that follow set down how that standing is entered, kept, and closed.

\begin{clauses}
\item Majority is reached upon a citizen's sixteenth naming-day.
\item A citizen not yet of majority is a ward of the household or floor to which they are assigned, and holds every duty this Pact lays upon a citizen of their years, but not the standing to swear before Judicial, to enter the lottery under Article 3, to speak in Assembly under Article 20, or to hold office.
\item No citizen may be stripped of citizenship for any cause. A citizen may be sent to cleaning under Article 18, or otherwise punished under Article 17, but remains a citizen, named in the rolls, until death is recorded under Section 5.
\end{clauses}

\section{The Naming of Children}

A child born without the lottery's grant under Article 3 is not thereby refused a name or entry into the rolls, but the citizens responsible for that birth are answerable under Article 17.

\begin{clauses}
\item Every child shall be named and entered into the rolls within ten days of birth, in a naming recorded by Judicial and attested by the delivering physician or midwife.
\item A child's naming-day is reckoned from the day of naming, not the day of birth, and every observance this Pact ties to a naming-day --- schooling under Article 15, majority under this Article, testing and shadowing under Article 13 --- is reckoned from it.
\item A child not named within ten days of birth shall be named regardless, upon inquiry by Judicial, and no delay in naming shortens or lengthens any citizen's later years of majority, testing, or retirement.
\end{clauses}

\section{The Citizen's Oath}

Upon a citizen's sixteenth naming-day, that citizen shall present before Judicial and swear the citizen's oath: to keep this Pact, to bear the labor given, and to hold the silo's peace above their own --- the same oath every citizen before them has sworn, without exception, since the first citizen stood where they now stand, and the same oath every citizen after them will swear in turn.

\begin{clauses}
\item A citizen who cannot present for the oath by reason of illness, distance of floor, or the press of labor may swear it before a deputy of Judicial sent for that purpose, but no citizen's majority is delayed beyond one turn of the season for want of the ceremony alone.
\item The oath once sworn is not sworn again, whatever office a citizen later holds or loses.
\end{clauses}

\section{The Keeping of the Rolls}

Judicial shall keep, for every living citizen, a roll recording that citizen's name, naming-day, floor and household of assignment, and trade, and shall issue to every citizen of majority an occupant card bearing the same, to be shown when Judicial, the Sheriff, Housing, or Supply require it.

\begin{clauses}
\item Housing, Supply, and every Head of Department shall report to Judicial, within ten days, any change to a citizen's floor, household, or trade, that the rolls may be kept current; Judicial is not charged to seek out such changes on its own account, and relies upon these reports being made in good faith.
\item The full count of the rolls, and no other number, is the silo's population for every purpose of this Pact, including the granting of a lottery under Article 3 and the setting of the ration under Article 11.
\item A citizen who believes the rolls to record them in error may petition Judicial for correction; Judicial's finding on such a petition is final.
\end{clauses}

\section{Death and the Closing of a Name}

A citizen's name is entered in the rolls at naming and closed in the rolls at death, and it is this closing, and no other event, that opens a place in the lottery under Article 3.

\begin{clauses}
\item A citizen's death shall be recorded by Judicial upon the attestation of a physician under Article 19, or, where a citizen is sent to cleaning under Article 18, upon the word of the Department of Information Technology that the sensors were kept.
\item No citizen may be recorded as dead while living, nor as living while dead, and a citizen found to have made a false attestation under this section is treated among the gravest offenses under Article 17.
\item A death suspected by the attending physician to have come by the citizen's own hand, and not sought and granted under Article 18, shall not be closed upon the physician's attestation alone, but shall be carried to Judicial for a finding under Article 6 before the name is closed.
\end{clauses}

\section{The Ceremony of the Oath}

The oath under Section 3 is sworn once, aloud, before witnesses, and in fixed words; this section sets down the manner of it.

\begin{clauses}
\item A citizen presents for the oath on the citizen's sixteenth naming-day, or on the first posted sitting of Judicial under Article 6 thereafter, in the hall of Judicial, or before a deputy of Judicial sent to the citizen's floor where Section 3 allows; the citizen's household attends where it will, and the citizen's teacher under Article 15 where the teacher can climb.
\item Before the oath is sworn, the Clerk of Judicial reads aloud the entry of the citizen's final examination under Article 15, Section 10, and the citizen's name, naming-day, floor, household, and the trade or shadowing the citizen enters under Article 13, and asks the citizen whether each is true; a thing the citizen says is not true is entered as disputed and carried to the rolls for correction under Section 4, and the oath is sworn regardless.
\item The oath is sworn in these words and no others, the citizen standing, the Clerk or deputy reading each line and the citizen speaking it after: \emph{I am} [the citizen's name], \emph{of} [the citizen's floor and household]. \emph{I will keep this Pact. I will bear the labor given me. I will hold the silo's peace above my own. I swear it before Judicial, as every citizen before me has sworn, and as every citizen after me will swear.}
\item The Clerk enters the oath in the rolls with the day, the place, and the witnesses, and gives the citizen the occupant card under Article 21 in the same hour; the citizen is of majority from the entry, and every standing this Pact gives a citizen of majority is the citizen's from that hour.
\item A citizen who will not speak the oath is not of majority and is not forced; the Clerk enters the refusal, the citizen remains a ward of the household under Section 1, and the citizen may present again at any later sitting. A ward of eighteen naming-days who has not sworn is carried before the Judge, who hears the ward and the household and rules whether the ward is fit to swear, or is to remain a ward of the silo under Article 3 in the Infirmary's keeping under Article 19.
\end{clauses}

\section{The Reconciling of the Rolls}

The rolls are fed by many hands under Section 4, and once in every year the hands are compared with one another.

\begin{clauses}
\item At the turn of the season the Judge names, Housing under Article 14, Supply under Article 11, the Infirmary under Article 19, the schools under Article 15, and the Head of every Department deliver to the Clerk of Judicial their own rolls of the citizens in their keeping --- the households and their dwellings, the drawings of the ration by name, the births and deaths attended, the children enrolled, the citizens assigned to each trade --- as those rolls stand on that day.
\item The Clerk and the archivist compare each delivered roll against the rolls of citizenship, name by name, and enter every difference found: a citizen in one roll and not another, a floor or household or trade entered differently, a name closed in one and open in another.
\item A difference is resolved upon the report of the office whose roll is the latest in its entry, and the rolls of citizenship corrected to it under Section 4; Judicial is not charged to seek out the citizen concerned, nor to inquire beyond the reports, and relies upon them as Section 4 provides. A difference the reports themselves cannot resolve is entered as open, and the citizen's occupant card under Article 21 is called in and compared.
\item The count of differences found, and of those left open, is entered in the archive and reported by the Judge to the Mayor and read at the next sitting of the Assembly under Article 20, without the names.
\end{clauses}

\section{The Count of the Rolls}

The number this Pact turns upon is a number, and is read as one.

\begin{clauses}
\item At every turn of the season, the Clerk of Judicial counts the names in the rolls not closed under Section 5, and enters the count with the day, the names entered since the last count, and the names closed; the count is delivered the same day to the Mayor and to the Head of Supply, and the ration under Article 11 and the lottery under Article 3 are reckoned from it and from no other number until the next.
\item A name neither closed nor confirmed by any report under Section 4 within the year --- a citizen no office has entered in any roll delivered under Section 7 --- stands open in the count, and is counted, until Judicial has inquired of the household and the floor and either confirmed or closed it; a name so open is reported to the Sheriff, who seeks the citizen.
\item The count is read aloud at the ordinary sitting of the Assembly under Article 20 with the reckoning of the stores, and any citizen may ask the Clerk the count at the door of Judicial on any posted day.
\end{clauses}
